\subsection{Edge Deployment}
\label{subsec:edge-deployment}

VisionServe targets resource-constrained edge scenarios without compromising inference
performance. The execution-provider fallback chain---TensorRT $\to$ CUDA $\to$ CoreML
$\to$ DirectML $\to$ OpenVINO $\to$ CPU---allows the same binary to exploit the best
available accelerator on each platform: NVIDIA Jetson modules via the TensorRT EP,
Raspberry Pi-class ARM boards via the CPU EP, and Apple Silicon Macs via the CoreML EP.
Because the server is a single statically-linked Go binary with no Python runtime
dependency, it eliminates the roughly 200\,MB overhead that PyTorch or TensorFlow
deployments impose even before any model weights are loaded. This footprint reduction
is significant on boards with 4--8\,GB of unified memory.

As a concrete example, rf-detr-nano runs on a Jetson Orin with the TensorRT EP; on-device
latency and energy on the Jetson are to be reported once measured on the device (we do not
relabel the workstation figures here). The same model falls back gracefully to the CUDA EP
on a workstation GPU, and to the CPU EP on a plain ARM board, with no configuration
change required from the operator.

\subsection{Text-Prompted Segmentation via Grounded-SAM}
\label{subsec:grounded-sam}

Grounded-SAM~\cite{liu2023groundingdino,zhang2023mobilesam} combines open-vocabulary
detection with prompted segmentation in a single HTTP request. The client submits a
natural-language prompt (e.g., ``red traffic cone'') together with the image; internally,
the \texttt{PipelineModel} implementation routes the request through two ONNX sessions
managed by the \texttt{lifecycle.Manager}: GroundingDINO~\cite{liu2023groundingdino}
first localises every object matching the text query and returns bounding boxes, which
are then passed as box prompts to MobileSAM~\cite{zhang2023mobilesam}, whose decoder
produces a binary mask for each detection. The two chained calls are invisible to the
caller---the server returns a single unified \texttt{Result} containing both the
detection bounding boxes and the corresponding RLE-encoded masks.

A minimal invocation illustrates the simplicity of the client interface:

\begin{verbatim}
curl -X POST http://localhost:11435/infer \
  -F "model=grounded-sam" \
  -F "prompt=red traffic cone" \
  -F "image=@scene.jpg"
\end{verbatim}

The response carries a \texttt{detections} array with bounding boxes in original image
coordinates and a \texttt{masks} array encoded as column-major RLE, ready for
downstream use without any client-side chaining logic.

\subsection{Depth-Aware Detection}
\label{subsec:depth-aware}

VisionServe's Python client exposes a \texttt{get\_depth\_at\_detection()} utility that
fuses monocular depth estimation with object detection in two independent server calls,
enabling spatial reasoning without a stereo camera or LiDAR sensor.

\begin{enumerate}
    \item Send the image to the MiDaS~\cite{ranftl2021midas} model endpoint to obtain a
          per-pixel relative depth map (\texttt{DepthResult}).
    \item Send the same image to an object detector (e.g., rf-detr) to obtain bounding
          boxes (\texttt{DetectionResult}).
    \item Call
          \texttt{depths = get\_depth\_at\_detection(depth\_result, det\_result)},
          which samples the depth map at the centroid of each bounding box and returns a
          per-detection depth value aligned with the detection list.
\end{enumerate}

The JavaScript/TypeScript client provides an equivalent \texttt{getDepthAtDetection()}
function with the same semantics. Practical applications include warehouse robotics
(estimating pick distance), autonomous navigation (obstacle ranging), and augmented
reality scene understanding (anchoring virtual objects at physically plausible depths).

\subsection{Open-Vocabulary and Domain-Specific Lightweight Models}
\label{subsec:open-vocab-domain}

Beyond general-purpose detection and segmentation, VisionServe supports
domain-specific lightweight models that conform to the same unified schema. PaddleOCR
\cite{du2020ppocr} implements a two-stage text detection and recognition pipeline
entirely within the \texttt{PipelineModel} interface: a detection stage localises text
regions as bounding boxes, and a recognition stage transcribes each region; the results
are returned as a standard \texttt{Result} with no schema changes required on the client.

SCRFD~\cite{guo2021scrfd} provides high-speed face detection at 45\,ms median latency
and approximately 420\,MB VRAM. Because it implements the plain \texttt{Model}
interface, adding face detection to an application requires only selecting the
\texttt{scrfd} model identifier---no task-specific client code is needed. Both examples
demonstrate that the architecture accommodates heterogeneous model families while
maintaining a single, predictable API surface.
